\begin{abstract}
The Indian mango (Mangifera indica L., locally known as Katchamitha) is an economically 
significant fruit crop in the Philippines, prized for its accessibility and distinct 
flavor. However, the manual assessment of its morphological properties, such as 
volume, weight, and density, is labor-intensive, subjective, and often destructive. 
This study proposes a novel, non-destructive vision-based framework for estimating 
the morphological properties of Katchamitha mangoes by integrating Monocular Depth 
Estimation (MDE) with instance segmentation. Due to the deceptively green coloration 
of Indian mangoes even when ripe, traditional RGB-based segmentation often 
struggles with green-to-green differentiation. To overcome this, a custom-trained 
YOLOv11n-seg model is employed for initial segmentation (achieving box and mask 
precision of 99.9\% and mAP50 of 99.5\%), and its output is refined using the depth 
map generated by Depth Anything V2. The study extracts geometric and color features,
including dimension ratios, solidity, curvature, and CIELab color values, from a 
dataset of 45 Katchamitha mango samples containing 243 raw multi-view images 
(expanded to 640 images via flip and rotation augmentations). 

We evaluate two distinct architectures: (1) an uncalibrated pipeline using raw 
monocular features, and (2) a calibrated pipeline using ArUco reference markers to 
resolve MDE scale ambiguity. In both, a unified MobileNetV3-Small CNN orientation 
classifier first determines the mango's pose (side, top, bottom, front, or back) to 
route the image crop to specialized Layer 1 Random Forest models for external 
morphology (height, width, thickness, volume, weight) and ripeness. A Layer 2 model 
then maps these outputs to internal seed dimensions via allometric relationships. 

Results show that the unified CNN orientation classifier achieves \textbf{97.32\%}
accuracy on unseen test subjects, a significant improvement over the previous
feature-based approach. The uncalibrated pipeline, however, fails to predict
physical dimensions (negative $R^2$ scores) due to depth scaling ambiguity. In
contrast, the ArUco-calibrated pipeline successfully resolves this ambiguity,
achieving $R^2$ scores of 0.83--0.87 for volumetric estimation, 0.78--0.86 for
weight prediction, and 0.90--0.94 for internal seed properties. This approach
offers a robust and cost-effective alternative to existing 2D and destructive
methods, paving the way for improved automated post-harvest grading and quality
assessment.
\end{abstract}

\begin{IEEEkeywords}
Monocular Depth Estimation, Instance Segmentation, Morphological Properties, Mango
Grading, Computer Vision, YOLOv11, Depth Anything V2, MobileNetV3, Orientation
Classification
\end{IEEEkeywords}
